\section*{Abstract}
		This work presents a comprehensive analysis of temperature units in natural units ($\hbar = c = k_B = 1$) within the framework of T0 theory. In the static $\xi$ universe, no expanding spacetime is needed. The derivations use the universal constant $\xi = \frac{4}{3} \times 10^{-4}$ as the only parameter (plus the integer settings named in the corpus; SI values serve for conversion) and respect the fundamental time-energy duality. The document includes CMB calculations within the framework of T0 theory, addresses fundamental questions regarding redshift mechanisms, primordial perturbations, and the resolution of cosmological tensions. The theory interprets the CMB without inflation and without a decoupling at any cosmological $z$ -- through stationary thermalization of the $\xi$ field in an infinitely old universe -- derives primordial perturbations from T-field quantum fluctuations, and addresses the Hubble tension (see Document~026 for the description calibrated to $H_0$, with $H_0 \approx 66.2\,$km/s/Mpc; the cosmic sector is calibrated, not derived from $\xi$).

	\medskip\noindent\textbf{Status markers.} Statements marked \textbf{[S]} are \emph{speculative}: a hint or conjecture that has not yet been derived or numerically checked in the corpus. The present corpus deliberately leaves the cosmic sector ($H_0$, $\Lambda$, dark energy, origin of the CMB) aside, because the Standard Model does not derive it either (P39); the cosmological interpretations of this document (static, eternal universe, origin of the CMB, comparison with $\Lambda$CDM) are to be read in this sense.

	
	
	\section{Introduction: T0 Theory in Natural Units}
	
	\subsection{Natural Units as Foundation}
	
	\begin{important}
		This entire work uses exclusively natural units with $\hbar = c = k_B = 1$. All quantities have energy dimensions: $[L] = [T] = [E^{-1}]$, $[M] = [T_{\text{temp}}] = [E]$.
	\end{important}
	
	The system of natural units represents a fundamental simplification of physics by setting the universal constants $\hbar$ (reduced Planck constant), $c$ (speed of light), and $k_B$ (Boltzmann constant) to the value 1. This choice simplifies the formulas and makes the relations between quantities directly visible.
	
	In this system, all of physics reduces to a single fundamental dimension -- energy. All other physical quantities are expressed as powers of energy:
	\begin{align}
		\text{Length:} \quad [L] &= [E^{-1}] \quad \text{(energy}^{-1}\text{)} \\
		\text{Time:} \quad [T] &= [E^{-1}] \quad \text{(energy}^{-1}\text{)} \\
		\text{Mass:} \quad [M] &= [E] \quad \text{(energy)} \\
		\text{Temperature:} \quad [T_{\text{temp}}] &= [E] \quad \text{(energy)}
	\end{align}
	
	This dimensional reduction makes relations between quantities transparent. In natural units, for example, Einstein's famous formula $E = mc^2$ becomes the trivial statement $E = m$, since both energy and mass have the same dimension.
	
	\textbf{Unit conversion (for reference):}
	For readers familiar with SI units, the following conversion factors apply:
	\begin{itemize}
		\item $\hbar = 1.055 \times 10^{-34}$ J$\cdot$s $\rightarrow 1$ (nat. units)
		\item $c = 2.998 \times 10^8$ m/s $\rightarrow 1$ (nat. units)  
		\item $k_B = 1.381 \times 10^{-23}$ J/K $\rightarrow 1$ (nat. units)
	\end{itemize}
	
	\subsection{The Universal $\xi$ Constant}
	
	\begin{important}[Key statement]
		T0 theory uses a single geometric constant $\xi = \frac{4}{3} \times 10^{-4}$ as its basic parameter; quantities from quarks to cosmic structures are derived from it. Cosmologically, this document describes a static, eternally existing cosmos without a Big Bang \textbf{[S]}. The factor $\frac{4}{3}$ originates from the fundamental geometric ratio between sphere volume and tetrahedron volume in three-dimensional space.
	\end{important}
	
	The heart of T0 theory is formed by a universal dimensionless constant, which we denote by the Greek letter $\xi$ (Xi). This constant was originally derived purely geometrically from the fundamental T0 field equations, as shown in the established T0 theory \cite{T0Theory}.
	
	The fundamental T0 theory is based on the universal dimensionless constant:
	\begin{equation}
		\xi = \frac{4}{3} \times 10^{-4} \quad \text{(dimensionless, exact geometric value)}
	\end{equation}
	
	\textbf{Geometric derivation from T0 field equations:} The value of $\xi$ follows directly from the geometric structure of the T0 field equations of the universal energy field $E_{\text{field}}(x,t)$. The fundamental T0 equation $\square E_{\text{field}} = 0$ in conjunction with three-dimensional spatial geometry leads to:
\begin{itemize}
	\item The geometric factor $\frac{4}{3}$ from three-dimensional spatial geometry
	\item The scale ratio $10^{-4}$ from the fractal dimension
	\item For the complete derivation see 041\_parameterherleitung\_De.pdf 
\end{itemize}
	
	\textbf{Comparison with experiment:} The deviation of the anomalous magnetic moment of the muon from the Standard Model is no longer significant as of 2025; it therefore does not serve as evidence for $\xi$ (Doc.~382). The corpus's $g-2$ line rests on ratios between the leptons; its test statement is the tau prediction $a_\tau \approx 1.281\times10^{-3}$, anchored to measured quantities (Doc.~382).
	
	This constant determines in T0 theory a wide range of physical phenomena:
	\begin{itemize}
		\item \textbf{Particle physics}: All elementary particle masses result from geometric quantum numbers $(n,l,j,r,p)$ scaled with $\xi$
		\item \textbf{Field theory}: Characteristic energy scales of all interactions follow from $\xi$ field dynamics
		\item \textbf{Gravitation}: The gravitational constant in natural units $G_{\text{nat}} = 2.61 \times 10^{-70}$ is a direct function of $\xi$
		\item \textbf{Cosmology} \textbf{[S]}: Thermodynamic equilibrium in the static, infinitely old universe is maintained through $\xi$ field cycles
	\end{itemize}
	
	\textbf{Symbol explanation:}
	\begin{itemize}
		\item $\xi$ (Xi): Universal dimensionless constant of T0 theory
		\item $E_\xi$: Characteristic energy scale, defined as $E_\xi = 1/\xi$
		\item $T_\xi$: Characteristic temperature, equal to $E_\xi$ in natural units
		\item $L_\xi$: Characteristic length scale of the $\xi$ field
		\item $G_{\text{nat}}$: Gravitational constant in natural units
		\item $\alpha_{\text{EM}}$: Electromagnetic coupling (= 1 in natural units by definition)
		\item $\beta$: Dimensionless parameter $\beta = r_0/r = 2GE/r$
		\item $\omega$: Photon energy (dimension $[E]$ in natural units)
	\end{itemize}
	
	\textbf{Coupling constants in natural units:}
	\begin{align}
		\alpha_{\text{EM}} &= 1 \quad \text{(by definition in natural units)} \\
		\alpha_G &= \xi^2 = \left(\frac{4}{3} \times 10^{-4}\right)^2 = 1.78 \times 10^{-8} \\
		\alpha_W &= \xi^{1/2} = \left(\frac{4}{3} \times 10^{-4}\right)^{1/2} = 1.15 \times 10^{-2} \\
		\alpha_S &= \xi^{-1/3} = \left(\frac{4}{3} \times 10^{-4}\right)^{-1/3} = 19.57
	\end{align}
	For comparison: the measured value is $\alpha_s(M_Z) \approx 0.118$.
	
	\textbf{Important clarification on units:}
	In this entire document, we work exclusively in natural units with $\hbar = c = k_B = 1$. This means:
	\begin{itemize}
		\item The electromagnetic coupling constant is $\alpha_{\text{EM}} = 1$ by definition (not 1/137 as in SI units)
		\item All other coupling constants are expressed relative to $\alpha_{\text{EM}} = 1$
		\item Energy, mass, and temperature have the same dimension
		\item Length and time have dimension energy$^{-1}$
	\end{itemize}
	
	\textbf{Dimensional consistency:} Since $\xi$ is purely dimensionless, it has the same value in all unit systems. In T0 theory it characterizes the geometry of the spacetime continuum and is treated as a natural constant, comparable to the fine-structure constant.
	
	\subsection{Time-Energy Duality and Static Universe}
	
	\begin{important}
		\textbf{[S]} This document interprets the universe as eternally existing, without a temporal beginning. The reason why there is no singular beginning follows from the time-energy duality $T\cdot E = 1$ with the lower time bound $T_0 = \xi\,t_P$ and the boundaryless torus (Doc.~306). The cosmological interpretation remains an interpretation, not a proof; the cosmic sector is deliberately left aside in the present corpus (P39).
	\end{important}
	
	From the time-mass duality $T\cdot m = 1$ (Doc.~003/010), the time-energy duality follows in natural units:
	\begin{equation}
		T \cdot E = 1 .
	\end{equation}
	
	Hence $E\to\infty \Leftrightarrow T\to 0$. In T0 theory, however, time is bounded from below by $T_0 = \xi\,t_P$, the time dressing of the minimal length $L_0 = \xi\,\ell_P$; correspondingly, energy is bounded from above by $E_\text{max} = E_P/\xi$, with exact closure $T_0\,E_\text{max} = 1$. An infinite energy density is thus structurally excluded. In addition, the torus is compact and boundaryless, $\partial T^4 = \emptyset$: the question of a beginning at $t=0$ presupposes a boundary, which a compact torus does not have (Doc.~306).
	
	Heisenberg's energy-time uncertainty relation $\Delta E \times \Delta t \geq 1/2$ does not carry this conclusion: a finite $\Delta t$ yields a \emph{finite} bound on $\Delta E$, not $\Delta E\to\infty$, and the $\Delta t$ of the relation is the time of change of an observable, not a cosmic age (Doc.~306).
	
	\textbf{Conclusion in this document [S]:} There is no singular beginning; this document therefore interprets the universe as eternally existing. This leads to the static T0 universe, which possesses the following properties:
	
	The T0 universe is therefore:
	\begin{itemize}
		\item \textbf{Static}: No expanding space -- the spacetime metric is time-independent
		\item \textbf{Eternal}: Without temporal beginning or end
		\item \textbf{Thermodynamically balanced}: A dynamic equilibrium is maintained through $\xi$ field cycles
		\item \textbf{Structurally stable}: Continuous formation and renewal of matter and structures
	\end{itemize}
	
	\textbf{Unit check of the time-energy duality:}
	\begin{align}
		[T] \times [E] &= [E^{-1}] \times [E] = [E^0] = \text{dimensionless} \\
		[1] &= \text{dimensionless} \quad \text{(satisfied)}
	\end{align}
	
	\section{$\xi$ Field and Characteristic Energy Scales}
	
	\subsection{$\xi$ Field as Universal Energy Mediator}
	
	\begin{formula}
		The universal constant $\xi = \frac{4}{3} \times 10^{-4}$ defines the fundamental energy scale of T0 theory:
		\begin{equation}
			E_\xi = \frac{1}{\xi} = \frac{1}{\frac{4}{3} \times 10^{-4}} = \frac{3}{4} \times 10^4 = 7500
		\end{equation}
		(all quantities in natural units)
	\end{formula}
	
	The $\xi$ field represents the fundamental energy field of the universe, from which all other fields and interactions emerge. Its characteristic energy scale $E_\xi$ results as the reciprocal of the dimensionless constant $\xi$.
	
	\textbf{Unit check for $E_\xi$:}
	\begin{align}
		[E_\xi] &= \left[\frac{1}{\xi}\right] = \frac{[E^0]}{[E^0]} = [E^0] = \text{dimensionless}
	\end{align}
	
	In natural units, dimensionless is equivalent to an energy unit, since all quantities are reduced to powers of energy. Therefore $[E_\xi] = [E]$.
	
	This characteristic energy directly corresponds to a characteristic temperature in natural units, since energy and temperature have the same dimension:
	\begin{equation}
		T_\xi = E_\xi = \frac{3}{4} \times 10^4 = 7500 \quad \text{(nat. units)}
	\end{equation}
	
	\textbf{Unit check for $T_\xi$:}
	\begin{align}
		[T_\xi] = [E_\xi] = [E] = [T_{\text{temp}}] \quad \text{(satisfied)}
	\end{align}
	
	\textbf{Physical interpretation:} The energy scale $E_\xi = 7500$ in natural units corresponds to an extremely high temperature characteristic of the fundamental processes of the $\xi$ field. This energy lies far above all known particle energies.
	
	\subsection{Characteristic $\xi$ Length Scale}
	
	The $\xi$ field also defines a characteristic length scale:
	\begin{equation}
		L_\xi = \frac{1}{E_\xi} = \frac{1}{7500} \approx 1.33 \times 10^{-4} \quad \text{(nat. units)}
	\end{equation}
	
	This length scale plays a fundamental role in the geometric structure of spacetime and appears in various physical phenomena.
	
	\section{CMB in T0 Theory: Static $\xi$ Universe}
	
	\subsection{CMB without Big Bang}
	
	\begin{important}[Cosmic sector]
		\textbf{[S]} According to the argument above, time-energy duality rules out a Big Bang; the CMB background radiation would then have a different origin than z=1100 decoupling. This classification is speculative (P39).
	\end{important}
	
	T0 theory interprets the cosmic microwave background radiation through $\xi$ field mechanisms \textbf{[S]}:
	
	\subsubsection{1. $\xi$ Field Quantum Fluctuations}
	The omnipresent $\xi$ field generates vacuum fluctuations with a characteristic energy scale. The exact dependence is derived from the measured ratio $T_{\text{CMB}}/E_\xi \approx \xi^2$.
	
	\subsubsection{2. Stationary Thermalization}
	In an infinitely old universe, the background radiation reaches thermodynamic equilibrium at the characteristic $\xi$ temperature.
	
	\begin{sibox}
		\textbf{CMB measurements (for reference only, in SI units):}
		\begin{itemize}
			\item Vacuum energy density: $\rho_{\text{vacuum}} = 4.17 \times 10^{-14}$ J/m$^3$
			\item Radiation power: $j = 3.13 \times 10^{-6}$ W/m$^2$
			\item Temperature: $T = 2.7255$ K
		\end{itemize}
	\end{sibox}
	
	\subsection{The Already Established $\xi$ Geometry}
	
	\begin{important}
		T0 theory had already established a fundamental length scale before the CMB analysis was performed. The CMB energy density is compared below with this already existing $\xi$ geometric structure.
	\end{important}
	
	From the original T0 theory formulation followed:
	
	\textbf{Characteristic mass:}
	\begin{equation}
		m_{\text{char}} = \frac{\xi}{2\sqrt{G_{\text{nat}}}} \approx 4.13 \times 10^{30} \quad \text{(nat. units)}
	\end{equation}
	
	\textbf{Universal scaling rule:}
	\begin{equation}
		\text{Factor} = 2.42 \times 10^{-31} \cdot m \quad \text{(for arbitrary mass } m \text{ in nat. units)}
	\end{equation}
	
	\textbf{Gravitational constant derived from $\xi$:}
	\begin{equation}
		G_{\text{nat}} = 2.61 \times 10^{-70} \quad \text{(nat. units)}
	\end{equation}
	% ================== COMPLETE CMB SECTION FROM 039_Zwei-Dipole-CMB_De.pdf ==================
	
	\section{The T0 Theory Framework for CMB}
	\label{sec:t0_framework}
	
	T0 theory proposes an extension of standard cosmology through the introduction of an intrinsic time field $\Tfield$ that couples to all matter and radiation. This theory arose from dissatisfaction with quantum mechanical nonlocality and the need for a deterministic framework that preserves causality while simultaneously explaining observed correlations.
	
	\subsection{Fundamental Postulates}
	
	T0 theory is based on three fundamental postulates:
	
	\begin{enumerate}
		\item \textbf{Time-mass duality}: The fundamental relation
		\begin{equation}
			\Tfield \cdot m(x) = 1
			\label{eq:time_mass_duality}
		\end{equation}
		
		\item \textbf{Universal coupling parameter}: A single parameter
		\begin{equation}
			\xipar = \frac{4}{3} \times 10^{-4}, \qquad \frac{\lambda_h^2 v^2}{16\pi^3 m_h^2} \approx 1.30 \times 10^{-4}
			\label{eq:xi_definition}
		\end{equation}
		checked against Higgs physics to about $2.3\,\%$, governs all T-field interactions. With the SM self-coupling $\lambda_h = m_h^2/(2v^2) \approx 0.129$ one has $\lambda_h^2 v^2/(16\pi^3 m_h^2) = m_h^2/(64\pi^3 v^2) \approx 1.30 \times 10^{-4}$, about $2.3\,\%$ below $4/30000$ (PDG 2024: $m_h = 125.20$~GeV, $v = 246.22$~GeV); this is a consistency check of $\xi$, not an exact derivation (origin in Doc.~385). The factor $\frac{4}{3}$ ultimately originates from the fundamental geometric ratio between sphere volume and tetrahedron volume in three-dimensional space.
		
		\item \textbf{Modified Robertson-Walker metric}:
		\begin{equation}
			ds^2 = -c^2dt^2[1 + 2\xipar\ln(a)] + a^2(t)[1 - 2\xipar\ln(a)]d\vec{x}^2
			\label{eq:modified_metric}
		\end{equation}
	\end{enumerate}
	
	\section{Power Spectrum Calculations}
	\label{sec:power_spectra}
	
	\subsection{Temperature Power Spectrum}
	
	The CMB temperature power spectrum is:
	
	\begin{equation}
		C_\ell^{TT} = \frac{2}{\pi}\int_0^\infty k^2 dk \, \mathcal{P}_\Psi(k) |\Theta_\ell(k,\eta_0)|^2 \times \left(1 + \xipar f_\ell(k)\right)
		\label{eq:cl_tt}
	\end{equation}
	
	where:
	\begin{equation}
		f_\ell(k) = \ln^2\left(\frac{k}{k_*}\right) - 2\ln\left(\frac{k}{k_*}\right)
	\end{equation}
	
	\subsection{E-Mode Polarization}
	
	\begin{equation}
		C_\ell^{EE} = \frac{2}{\pi}\int_0^\infty k^2 dk \, \mathcal{P}_\Psi(k) |E_\ell(k,\eta_0)|^2 \times \left(1 + \xipar g_\ell(k)\right)
	\end{equation}
	
	\subsection{Cross-Correlation}
	
	\begin{equation}
		C_\ell^{TE} = \frac{2}{\pi}\int_0^\infty k^2 dk \, \mathcal{P}_\Psi(k) \Theta_\ell(k,\eta_0) E_\ell^*(k,\eta_0) \times \left(1 + \xipar h_\ell(k)\right)
	\end{equation}
	
	\section{MCMC Analysis and Parameter Constraints}
	\label{sec:mcmc}
	
	\subsection{Bayesian Parameter Estimation}
	
	We perform an MCMC analysis with:
	
	\begin{equation}
		\mathcal{L} = -\frac{1}{2}\sum_{\ell} \frac{2\ell+1}{2} f_{\text{sky}} \left[\frac{C_\ell^{\text{obs}} - C_\ell^{\text{theory}}(\theta)}{\sigma_\ell}\right]^2
	\end{equation}
	
	\subsection{Results with Uncertainties}
	
	\begin{table}[htbp]
		\centering
		\caption{T0 parameter constraints (68\% CL)}
		\adjustbox{max width=\textwidth}{%
\begin{tabular}{lcc}
			\toprule
			Parameter & Best fit & Uncertainty \\
			\midrule
			$H_0$ [km/s/Mpc] & 66.2$^*$ & calibrated to $H_0$ \\
			$\Omega_b h^2$ & 0.02237 & $\pm 0.00015$ \\
			$\Omega_c h^2$ & 0.1200 & $\pm 0.0012$ \\
			$\tau$ & 0.054 & $\pm 0.007$ \\
			$n_s$ & 0.9649 & $\pm 0.0042$ \\
			$\ln(10^{10}A_s)$ & 3.044 & $\pm 0.014$ \\
			$\xipar$ & $\frac{4}{3} \times 10^{-4}$ & (geometric constant) \\
			\bottomrule
		\end{tabular}%
}
		\label{tab:parameters}
	\end{table}
{$^*$The $H_0$-calibrated T0 description $H_0^{\text{T0}} = E_H/\hbar \approx 66.2\,$km/s/Mpc ($-1.9\%$ deviation from Planck) as well as the proposed route to resolving the Hubble tension are centralized in Document~026 (Geometric Cosmology).}
	
	\subsection{$S_8$ Tension}
	
	The clustering amplitude is modified:
	
	\begin{equation}
		S_8^{T0} = S_8^{\Lambda\text{CDM}} \times (1 - 2\xipar) = 0.834 \times (1 - 2 \times \frac{4}{3} \times 10^{-4}) = 0.834 \times 0.99973 = 0.8338
	\end{equation}
	
	In the assessment of this document, this is compatible with weak lensing measurements \textbf{[S]}.
	
	\section{Experimental Predictions}
	\label{sec:predictions}
	
	\subsection{Testable Predictions}
	
	T0 theory makes several specific predictions; as cosmological statements they are speculative \textbf{[S]} (P39):
	
	\begin{enumerate}
		\item \textbf{Running of the spectral index}:
		\begin{equation}
			\frac{dn_s}{d\ln k} = -2\xipar = -2 \times \frac{4}{3} \times 10^{-4} = -2.67 \times 10^{-4}
		\end{equation}
		
		\item \textbf{Tensor-to-scalar ratio}:
		\begin{equation}
			r = 16\xipar = 16 \times \frac{4}{3} \times 10^{-4} = 0.00213 \pm 0.0004
		\end{equation}
		
		\item \textbf{Modified Silk damping}:
		\begin{equation}
			C_\ell^{TT} \propto \exp\left[-\left(\frac{\ell}{\ell_D}\right)^2\right] \times \left(1 + \xipar \left(\frac{\ell}{3000}\right)^2\right)
		\end{equation}
	\end{enumerate}
	
	A wavelength-dependent redshift $\Delta z \propto \ln(\lambda/\lambda_0)$ is not predicted: the fractal path lengthening is purely geometric and stretches all wavelengths by the same factor, so the redshift is achromatic (K6).
	
	\subsection{Observational Tests}
	
	\begin{table}[htbp]
		\centering
		\caption{T0 predictions vs observations}
		\adjustbox{max width=\textwidth}{%
\begin{tabular}{p{3cm}p{3cm}p{3cm}p{3cm}}
			\toprule
			Observable & T0 prediction & Current limit & Future sensitivity \\
			\midrule
			$dn_s/d\ln k$ & $-2.67 \times 10^{-4}$ & $< 0.01$ & $10^{-4}$ (CMB-S4) \\
			$r$ & $0.00213$ & $< 0.036$ & $0.001$ (LiteBIRD) \\
			$f_{NL}$ & $-3.5 \times 10^{-4}$ & $< 5$ & $0.1$ (CMB-S4) \\
			\bottomrule
		\end{tabular}%
}
	\end{table}
	
	\section{Comparison with $\Lambda$CDM}
	\label{sec:comparison}
	
	\subsection{$\chi^2$ Analysis}
	
	Comparison of model fits to Planck 2018 data:
	
	\begin{align}
		\chi^2_{\Lambda\text{CDM}} &= 1127.4 \\
		\chi^2_{T0} &= 1123.8 \\
		\Delta\chi^2 &= -3.6 \quad (1.9\sigma \text{ improvement})
	\end{align}
	
	The $\chi^2$ values themselves are not backed by any verification script.
	
	\subsection{Information Criteria}
	
	With the Akaike Information Criterion (AIC):
	
	\begin{equation}
		\Delta\text{AIC} = \Delta\chi^2 + 2\Delta N_{\text{params}} = -3.6 + 2 = -1.6
	\end{equation}
	
	In this comparison, the negative value slightly favors T0, despite the additional parameter \textbf{[S]}.
	
	\section{Modified Recombination History (\texorpdfstring{$\Lambda$CDM}{LambdaCDM} comparison frame)}

	\begin{important}
	\textbf{Comparison calculation, not an FFGFT-intrinsic quantity.} The static
	$\xi$ universe has no cosmological redshift $z$ in the expansion sense; what
	FFGFT calls redshift is fractal path-lengthening, not an expansion scale
	factor. The redshift variable $z$ used in this section is therefore a
	$\Lambda$CDM pipeline quantity, employed solely to compare numbers against
	the standard recombination calculation. It does not assert an FFGFT epoch of
	recombination; in FFGFT the CMB arises from stationary thermalization in an
	infinitely old universe (see the $\xi$-field mechanism above), not from a
	decoupling at any $z$.
	\end{important}

	Within this comparison frame, the standard recombination condition reads:
	\begin{equation}
		z_{\text{rec}} = \text{solution of } x_e(z) = 0.5
	\end{equation}

	The electron fraction evolves as:
	\begin{equation}
		x_e(z) = \frac{1}{1 + A(T) \exp[E_I/kT(z)]}
	\end{equation}

	where:
	\begin{align}
		T(z) &= T_0(1+z)[1 - \xi\ln(1+z)] \\
		A(T) &= \left(\frac{2\pi m_e kT}{h^2}\right)^{-3/2} 
		\frac{g_p g_e}{g_H} (1 + \xi h(T))
	\end{align}

	With the $\xi$ correction applied to the comparison, this yields
	$z_{\text{rec}} \approx 1089.5$, versus
	$z_{\text{rec}}^{\Lambda\text{CDM}} = 1089.9$ in the unmodified pipeline --
	a difference that exists only inside the $\Lambda$CDM comparison frame and is
	not an FFGFT prediction of a recombination epoch.

	% ================== END OF CMB SECTION ==================
	
	\section{CMB-Casimir Connection and $\xi$ Field Verification}
	\label{sec:cmb_casimir}
	
	\subsection{CMB Energy Density and $\xi$ Length Scale}
	
	\begin{important}[Thesis]
		\textbf{[S]} Thesis of this section: The measured CMB spectrum corresponds to the radiating energy density of the $\xi$ field vacuum; the vacuum itself radiates at its characteristic temperature.
	\end{important}
	
	The CMB energy density in natural units:
	\begin{equation}
		\rho_{\text{CMB}} = 2.0 \times 10^{-15}\ \text{eV}^4 \quad \text{(nat. units, dimension } [E^4] \text{)}
	\end{equation}
	
	With $T_{\text{CMB}} = 2.35 \times 10^{-4}$~eV one obtains $\rho_{\text{CMB}} = \frac{\pi^2}{15} T_{\text{CMB}}^4 = 2.0 \times 10^{-15}$~eV$^4$ and from it $L_\xi = (\xi/\rho_{\text{CMB}})^{1/4} = 508$~eV$^{-1} = 1.00 \times 10^{-4}$~m.
	
	The CMB temperature in natural units:
	\begin{equation}
		T_{\text{CMB}} = 2.35 \times 10^{-4} \quad \text{(nat. units)}
	\end{equation}
	
	This energy density defines a characteristic $\xi$ length scale:
	\begin{equation}
		L_\xi = \left(\frac{\xi}{\rho_{\text{CMB}}}\right)^{1/4}
	\end{equation}
	
	\begin{formula}
		Fundamental relation of CMB energy density:
		\begin{equation}
			\rho_{\text{CMB}} = \frac{\xi}{L_\xi^4} = \frac{\frac{4}{3} \times 10^{-4}}{L_\xi^4}
		\end{equation}
	\end{formula}
	
	\subsection{Casimir-CMB Ratio}
	
	The Casimir effect represents a direct manifestation of quantum vacuum fluctuations. In natural units, the Casimir energy density between two parallel plates with separation $d$ is (the formula with $240$ in the denominator describes the Casimir pressure, see the force calculation):
	
	\begin{equation}
		|\rho_{\text{Casimir}}| = \frac{\pi^2}{720 d^4} \quad \text{(nat. units)}
	\end{equation}
	
	At the characteristic $\xi$ length scale $L_\xi = 10^{-4}$ m, the ratio between Casimir and CMB energy densities provides a comparison value:
	
	\begin{equation}
		\frac{|\rho_{\text{Casimir}}|}{\rho_{\text{CMB}}} = \frac{\pi^2}{720 \xi} = \frac{\pi^2}{720 \times \frac{4}{3} \times 10^{-4}} = \frac{\pi^2 \times 10^4}{960} \approx 102.8
	\end{equation}
	
	\subsection{Detailed Calculations in SI Units}
	
	\textbf{Casimir energy density at plate separation} $d = L_\xi = 10^{-4}$ m:
	
	\begin{align}
		|\rho_{\text{Casimir}}| &= \frac{\hbar c \pi^2}{720 d^4} \\
		&= \frac{1.055 \times 10^{-34} \times 2.998 \times 10^8 \times \pi^2}{720 \times (10^{-4})^4} \\
		&= \frac{3.12 \times 10^{-25}}{7.2 \times 10^{-14}} \\
		&= 4.3 \times 10^{-12} \text{ J/m}^3
	\end{align}
	
	\textbf{CMB energy density in SI units:}
	\begin{equation}
		\rho_{\text{CMB}} = 4.17 \times 10^{-14} \text{ J/m}^3
	\end{equation}
	
	\textbf{Ratio in SI units:}
	\begin{equation}
		\frac{|\rho_{\text{Casimir}}|}{\rho_{\text{CMB}}} = \frac{4.3 \times 10^{-12}}{4.17 \times 10^{-14}} = 104
	\end{equation}
	
	\textbf{Theoretical prediction in natural units:}
	\begin{align}
		\frac{|\rho_{\text{Casimir}}|}{\rho_{\text{CMB}}} &= \frac{\pi^2 / (720 L_\xi^4)}{\xi / L_\xi^4} \\
		&= \frac{\pi^2}{720 \xi} = \frac{\pi^2}{720 \times \frac{4}{3} \times 10^{-4}} \\
		&= \frac{\pi^2 \times 3 \times 10^4}{720 \times 4} = \frac{\pi^2 \times 10^4}{960} \approx 102.8
	\end{align}
	
	\textbf{Agreement:} The ratio 104 calculated in SI units agrees with the value 102.8 in natural units to within 1.1\%. Both values follow from the same formula; their closeness follows from the definition of $L_\xi$ via $\rho_{\text{CMB}}$ and is not an independent comparison with a measurement; the ratio at $L_\xi$ is an identity.
	
	The deviation between the two values (102.8 and 104) is 1.1\% and corresponds to rounding $L_\xi$ to $10^{-4}$~m.
	
	\begin{important}
		At the characteristic $\xi$ length scale $L_\xi = 10^{-4}$ m, the CMB vacuum energy density and the Casimir energy density do not coincide (ratio about $10^2$); equality lies at a larger plate separation. The Casimir-CMB connection is a statement about scales \textbf{[S]}, not evidence for the $\xi$ field.
	\end{important}
	
	\subsection{Dimensionless $\xi$ Hierarchy and Independent Verification}
	
	\textbf{Critical question: Is this circular reasoning?}
	
	In the assessment of this document, there is no circular reasoning, because:
	
	\begin{enumerate}
		\item \textbf{Different theoretical and experimental sources:}
		\begin{itemize}
			\item $\xi$ constant: Purely geometrically derived from T0 field equations
			\item CMB data: Cosmic microwave measurements
			\item Casimir measurements: Laboratory vacuum experiments
		\end{itemize}
		
		\item \textbf{Temporal sequence of development:}
		\begin{itemize}
			\item T0 theory and $\xi$ derivation: Purely theoretical geometric derivation
			\item CMB prediction: Followed from the already established $\xi$ geometry
			\item Casimir comparison: identity via the definition of $L_\xi$, not an independent path
		\end{itemize}
		
		\item \textbf{Several comparison paths:}
		\begin{itemize}
			\item Geometric derivation $\to$ $\xi = \frac{4}{3} \times 10^{-4}$
			\item Higgs mechanism $\to$ $\frac{\lambda_h^2 v^2}{16\pi^3 m_h^2} \approx 1.30 \times 10^{-4}$, about $2.3\,\%$ below $\frac{4}{3} \times 10^{-4}$ (consistency check, see Eq.~\ref{eq:xi_definition})
			\item Lepton masses $\to$ $\xi = \frac{4}{3} \times 10^{-4}$
			\item CMB/Casimir ratio $\to$ identity via the definition of $L_\xi$, not a path of its own
		\end{itemize}
	\end{enumerate}
	
	\subsubsection{Detailed Energy Scale Ratios}
	
	The dimensionless ratio between CMB temperature and characteristic energy -- detailed calculation:
	
	\begin{align}
		\frac{T_{\text{CMB}}}{E_\xi} &= \frac{2.35 \times 10^{-4}}{\frac{3}{4} \times 10^4} \\
		&= \frac{2.35 \times 10^{-4} \times 4}{3 \times 10^4} \\
		&= \frac{9.4}{3 \times 10^8} \\
		&= \frac{9.4}{3} \times 10^{-8} \\
		&= 3.13 \times 10^{-8}
	\end{align}
	
	Theoretical prediction from $\xi$ geometry -- detailed steps:
	\begin{align}
		\xi^2 &= \left(\frac{4}{3} \times 10^{-4}\right)^2 \\
		&= \frac{16}{9} \times 10^{-8} \\
		&= 1.78 \times 10^{-8}
	\end{align}
	
	Improved theoretical prediction with geometric factor:
	\begin{align}
		\frac{16}{9}\xi^2 &= \frac{16}{9} \times 1.78 \times 10^{-8} \\
		&= 1.778 \times 1.78 \times 10^{-8} \\
		&= 3.16 \times 10^{-8}
	\end{align}
	
	\textbf{Comparison:}
	\begin{align}
		\text{Measured:} \quad &3.13 \times 10^{-8} \\
		\text{Theoretical:} \quad &3.16 \times 10^{-8} \\
		\text{Agreement:} \quad &\frac{3.13}{3.16} = 0.99 = 99\% \text{ (1\% deviation)}
	\end{align}
	
	Agreement to within 1\% for the relation
	\begin{equation}
		\boxed{\frac{T_{\text{CMB}}}{E_\xi} = \frac{16}{9}\xi^2}
	\end{equation}
	This agreement is unit-dependent: the relation gives a pure number that agrees with the measured value only in the energy unit eV. It is therefore a choice of scale, not a derivation; the forward derivation of the CMB temperature is open.
	
	\subsubsection{Length Scale Ratios}
	
	\begin{equation}
		\frac{\ell_{\xi}}{L_\xi} = \xi^{-1/4} = \left(\frac{3}{4}\right)^{1/4} \times 10
	\end{equation}
	
	\subsection{Consistency Verification of T0 Theory}
	
	\begin{important}[Consistency check]
		The $\xi$ constant derived from particle physics is compatible with the vacuum energy density measured from the CMB; reading this as a prediction belongs to the cosmic sector \textbf{[S]}.
	\end{important}
	
	The paths to the length scale compared:
	
	\begin{table}[htbp]
		\centering
		\caption{Consistency verification of the $\xi$ length scale}
		\adjustbox{max width=\textwidth}{%
\begin{tabular}{p{4cm}p{4cm}p{4cm}}
			\toprule
			\textbf{Derivation} & \textbf{Starting point} & \textbf{Result} \\
			\midrule
			$\xi$ geometry (bottom-up) & $\xi = \frac{4}{3} \times 10^{-4}$ from particles & $L_\xi = 1/E_\xi = 1.33 \times 10^{-4}$ (nat.\ units) \\
			CMB vacuum (top-down) & $\rho_{\text{CMB}}$ from measurement & $L_\xi = \left(\frac{\xi}{\rho_{\text{CMB}}}\right)^{1/4} = 1.0 \times 10^{-4}$ m \\
			Casimir effect & Casimir formula & Identity via $L_\xi$ from $\rho_{\text{CMB}}$ \\
			\midrule
			\textbf{Comparison} & \textbf{Two different quantities} & numerical coincidence \\
			\bottomrule
		\end{tabular}%
}
	\end{table}
	Two different quantities appear in this document under the name $L_\xi$. $L_\xi = 1/E_\xi = 1.33 \times 10^{-4}$ (section on the $\xi$ length scale) is a number in natural units: read in eV$^{-1}$ it corresponds, with $\hbar c = 1.973 \times 10^{-7}$~eV$\cdot$m, to a length of $2.6 \times 10^{-11}$~m, in Planck units ($E_{\text{Planck}} = 1$, as in the section on energy scales) to $2.2 \times 10^{-39}$~m; $L_\xi = (\xi/\rho_{\text{CMB}})^{1/4} = 508$~eV$^{-1}$ corresponds to $1.0 \times 10^{-4}$~m. The value $10^{-4}$~m in the \enquote{bottom-up} row therefore does not follow from $\xi$ alone but from $\rho_{\text{CMB}}$; the agreement $1.33 \times 10^{-4} \approx 10^{-4}$ is a numerical coincidence of two different quantities.
	
	\subsection{The $\xi$ Field as Universal Vacuum}
	
	\begin{formula}
		The $\xi$ field vacuum manifests in multiple phenomena:
		\begin{align}
			\text{Free vacuum (CMB):} \quad &\rho_{\text{CMB}} = \frac{\xi}{L_\xi^4} \\
			\text{Confined vacuum (Casimir):} \quad &|\rho_{\text{Casimir}}| = \frac{\pi^2}{720 d^4} \\
			\text{Ratio at } d = L_\xi: \quad &\frac{|\rho_{\text{Casimir}}|}{\rho_{\text{CMB}}} = \frac{\pi^2 \times 10^4}{960}
		\end{align}
	\end{formula}
	
	\begin{important}
		All $\xi$ relations consist of exact mathematical ratios:
		\begin{itemize}
			\item Fractions: $\frac{4}{3}$, $\frac{16}{9}$, $\frac{3}{4}$
			\item Powers of ten: $10^{-4}$, $10^4$
			\item Mathematical constants: $\pi^2$
		\end{itemize}
		No freely fitted decimal numbers occur; the ratios follow from $\xi$ geometry.
	\end{important}
	\section{Casimir Effect and $\xi$ Field Connection}
	
	\subsection{Modified Casimir Formula in T0 Theory}
	
	T0 theory offers an interpretation of the Casimir effect through the $\xi$ field:
	
	\begin{equation}
		|\rho_{\text{Casimir}}(d)| = \frac{\pi^2}{720 \xi} \rho_{\text{CMB}} \left(\frac{L_\xi}{d}\right)^4
	\end{equation}
	
	Substituting $\rho_{\text{CMB}} = \xi/L_\xi^4$ yields the standard formula:
	\begin{equation}
		|\rho_{\text{Casimir}}| = \frac{\pi^2}{720 d^4}
	\end{equation}
	
	This is compatible with the interpretation that the Casimir effect and the CMB are different manifestations of the same $\xi$ field vacuum \textbf{[S]}.
	
	\section{Structure Formation in the Static $\xi$ Universe}
	
	\subsection{Continuous Structure Development}
	
	In the static T0 universe \textbf{[S]}, structure formation takes place continuously without Big Bang constraints:
	
	\begin{equation}
		\frac{d\rho}{dt} = -\nabla \cdot (\rho \mathbf{v}) + S_\xi(\rho, T, \xi)
	\end{equation}
	
	where $S_\xi$ is the $\xi$ field source term for continuous matter/energy transformation.
	
	\subsection{$\xi$-Assisted Continuous Creation}
	
	In this picture \textbf{[S]}, the $\xi$ field enables continuous matter/energy transformation:
	
	\begin{align}
		\text{Quantum vacuum} &\xrightarrow{\xi} \text{Virtual particles} \\
		\text{Virtual particles} &\xrightarrow{\xi^2} \text{Real particles} \\
		\text{Real particles} &\xrightarrow{\xi^3} \text{Atomic nuclei} \\
		\text{Atomic nuclei} &\xrightarrow{\text{time}} \text{Stars, galaxies}
	\end{align}
	
	The energy balance is maintained by:
	\begin{equation}
		\rho_{\text{total}} = \rho_{\text{matter}} + \rho_{\xi\text{-field}} = \text{constant}
	\end{equation}
	
	\begin{important}
		\textbf{[S]} In this picture, energy is conserved through continuous transformation between matter and $\xi$ field energy, which allows an eternal existence without beginning or end.
	\end{important}
	
	\section{Unit Analysis of the $\xi$-Based Casimir Formula}
	
	This analysis examines the unit consistency of the modified Casimir formula within T0 theory, which introduces the dimensionless constant $\xi$ and the cosmic microwave background (CMB) energy density $\rho_{\text{CMB}}$. The goal is to verify consistency with the standard Casimir formula and to clarify the physical meaning of the new parameters $\xi$ and $L_\xi$. The analysis is performed in SI units, with each formula checked for dimensional correctness.
	
	\subsection{Standard Casimir Formula}
	The standard Casimir formula describes the energy density of the Casimir effect between two parallel, perfectly conducting plates in vacuum:
	\begin{equation}
		|\rho_{\text{Casimir}}| = \frac{\pi^2 \hbar c}{720 d^4}
	\end{equation}
	Here $\hbar$ is the reduced Planck constant, $c$ the speed of light, and $d$ the separation between the plates. The unit check yields:
	\begin{equation}
		\frac{[\hbar] \cdot [c]}{[d^4]} = \frac{(\text{J} \cdot \text{s}) \cdot (\text{m}/\text{s})}{\text{m}^4} = \frac{\text{J} \cdot \text{m}}{\text{m}^4} = \frac{\text{J}}{\text{m}^3}
	\end{equation}
	This corresponds to the unit of energy density and confirms the correctness of the formula.
	
	\textbf{Formula explanation:} The Casimir effect arises from quantum fluctuations of the electromagnetic field in vacuum. Only certain wavelengths fit between the plates, leading to a measurable energy density that scales as $d^{-4}$. The constant $\pi^2/720$ results from the summation over all allowed modes.
	
	\subsection{Definition of $\xi$ and CMB Energy Density}
	T0 theory introduces the dimensionless constant $\xi$, defined as:
	\begin{equation}
		\xi = \frac{4}{3} \times 10^{-4}
	\end{equation}
	This constant is dimensionless, confirmed by $[\xi] = [1]$. The CMB energy density is defined in natural units as:
	\begin{equation}
		\rho_{\text{CMB}} = \frac{\xi}{L_\xi^4}
	\end{equation}
	with the characteristic length scale $L_\xi = 10^{-4}$ m. In SI units, the CMB energy density is:
	\begin{equation}
		\rho_{\text{CMB}} = 4.17 \times 10^{-14} \text{ J}/\text{m}^3
	\end{equation}
	
	\textbf{Formula explanation:} The CMB energy density represents the energy of the cosmic microwave background radiation. In T0 theory, it is scaled by $\xi$ and $L_\xi$, where $L_\xi$ is a fundamental length scale possibly linked to cosmic phenomena. The unit analysis shows:
	\begin{equation}
		[\rho_{\text{CMB}}] = \frac{[\xi]}{[L_\xi^4]} = \frac{1}{\text{m}^4} = \text{E}^4 \text{ (in natural units)}
	\end{equation}
	In SI units, this gives J/m$^3$, which is consistent.
	
	\subsection{Conversion of the $\xi$ Relation to SI Units}
	T0 theory postulates a fundamental relation:
	\begin{equation}
		\hbar c \stackrel{!}{=} \frac{\rho_{\text{CMB}} L_\xi^4}{\xi}
	\end{equation}
	The unit analysis confirms:
	\begin{equation}
		\frac{[\rho_{\text{CMB}}] \cdot [L_\xi^4]}{[\xi]} = \left( \frac{\text{J}}{\text{m}^3} \right) \cdot \text{m}^4 \cdot 1 = \text{J} \cdot \text{m}
	\end{equation}
	This corresponds to the unit of $\hbar c$. Numerically we obtain:
	\begin{equation}
		\frac{\left( 4.17 \times 10^{-14} \right) \cdot \left( 10^{-4} \right)^4}{\frac{4}{3} \times 10^{-4}} = 3.13 \times 10^{-26} \text{ J} \cdot \text{m}
	\end{equation}
	
	The relation follows from the definition $\rho_{\text{CMB}} = \xi\hbar c/L_\xi^4$. Compared with $\hbar c = 3.16 \times 10^{-26}$ J$\cdot$m, this is $0.99$ times the value; the agreement to about $1\,\%$ follows from fixing $L_\xi$ via $\rho_{\text{CMB}}$.
	
	\textbf{Formula explanation:} This relation bridges quantum mechanics ($\hbar c$) with cosmic scales ($\rho_{\text{CMB}}$, $L_\xi$). The dimensionless constant $\xi$ acts as a scaling factor linking the CMB energy density with the fundamental length scale $L_\xi$.
	
	\subsection{Modified Casimir Formula}
	The modified Casimir formula is:
	\begin{equation}
		|\rho_{\text{Casimir}}(d)| = \frac{\pi^2}{720 \xi} \rho_{\text{CMB}} \left( \frac{L_\xi}{d} \right)^4
	\end{equation}
	The unit analysis yields:
	\begin{equation}
		\frac{[\rho_{\text{CMB}}] \cdot [L_\xi^4]}{[\xi] \cdot [d^4]} = \frac{\left( \frac{\text{J}}{\text{m}^3} \right) \cdot \text{m}^4}{1 \cdot \text{m}^4} = \frac{\text{J}}{\text{m}^3}
	\end{equation}
	This confirms the unit of energy density. Substituting $\rho_{\text{CMB}} = \xi \hbar c / L_\xi^4$ yields the standard Casimir formula:
	\begin{equation}
		|\rho_{\text{Casimir}}| = \frac{\pi^2}{720} \frac{\xi \hbar c}{L_\xi^4} \cdot \frac{L_\xi^4}{d^4} = \frac{\pi^2 \hbar c}{720 d^4}
	\end{equation}
	
	\textbf{Formula explanation:} The modified formula incorporates $\xi$ and $\rho_{\text{CMB}}$, linking the Casimir effect with cosmic parameters. Its consistency with the standard formula shows that T0 theory offers an alternative representation of the effect.
	
	\subsection{Force Calculation}
	The force per area is derived from the energy density:
	\begin{equation}
		\frac{F}{A} = -\frac{\partial}{\partial d} \left( |\rho_{\text{Casimir}}| \cdot d \right) = \frac{\pi^2}{240 \xi} \rho_{\text{CMB}} \left( \frac{L_\xi}{d} \right)^4
	\end{equation}
	
	With $\rho_{\text{CMB}} = \xi\hbar c/L_\xi^4$ this yields the literature formula of the Casimir pressure, $\pi^2\hbar c/(240\,d^4)$.
	The unit analysis shows:
	\begin{equation}
		\frac{[\rho_{\text{CMB}}] \cdot [L_\xi^4]}{[\xi] \cdot [d^4]} = \frac{\left( \frac{\text{J}}{\text{m}^3} \right) \cdot \text{m}^4}{1 \cdot \text{m}^4} = \frac{\text{J}}{\text{m}^3} = \frac{\text{N}}{\text{m}^2}
	\end{equation}
	This corresponds to the unit of pressure and confirms the correctness.
	
	\textbf{Formula explanation:} The force per area represents the measurable Casimir force arising from the change in energy density with plate separation. T0 theory scales this force with $\xi$ and $\rho_{\text{CMB}}$, enabling a cosmic interpretation.
	
	\subsection{Critical Assessment}
	T0 theory shows unit consistency; $\rho_{\text{CMB}} L_\xi^4/\xi$ matches $\hbar c$ to about $1\,\%$, which follows from fixing $L_\xi$ (no factor 16/9 occurs). It links the Casimir effect with cosmic vacuum energy through $\xi$ and $L_\xi$, where $L_\xi = 10^{-4}$ m serves as a fundamental length scale. This opens new physical interpretations connecting the Casimir effect with cosmological phenomena.
	
	\section{Dimensionless $\xi$ Hierarchy}
	
	\subsection{Complete Table of Dimensionless Ratios}
	
	All $\xi$ relations reduce to exact mathematical ratios:
	
	\begin{table}[htbp]
		\centering
		\caption{Dimensionless $\xi$ ratios in T0 theory}
		\adjustbox{max width=\textwidth}{%
\begin{tabular}{lcc}
			\toprule
			\textbf{Ratio} & \textbf{Expression} & \textbf{Value} \\
			\midrule
			Temperature ratio & $\frac{T_{\text{CMB}}}{E_\xi}$ & $3.13 \times 10^{-8}$ \\
			Relation (unit-dependent) & $\frac{16}{9}\xi^2$ & $3.16 \times 10^{-8}$ \\
			Length ratio & $\frac{\ell_{\xi}}{L_\xi}$ & $\xi^{-1/4}$ \\
			Casimir-CMB & $\frac{|\rho_{\text{Casimir}}|}{\rho_{\text{CMB}}}$ & $\frac{\pi^2 \times 10^4}{960}$ \\
			Gravitational coupling & $\alpha_G$ & $\xi^2 = 1.78 \times 10^{-8}$ \\
			Weak coupling & $\alpha_W$ & $\xi^{1/2} = 1.15 \times 10^{-2}$ \\
			Strong coupling & $\alpha_S$ & $\xi^{-1/3} = 19.57$ \\
			\bottomrule
		\end{tabular}%
}
	\end{table}
	
	\begin{important}
		All $\xi$ relations consist of exact mathematical ratios:
		\begin{itemize}
			\item Fractions: $\frac{4}{3}$, $\frac{3}{4}$, $\frac{16}{9}$
			\item Powers of ten: $10^{-4}$, $10^3$, $10^4$
			\item Mathematical constants: $\pi^2$
		\end{itemize}
		No freely fitted decimal numbers occur; the ratios follow from $\xi$ geometry.
	\end{important}
	
	\subsection{Parameter Reduction}
	
	\begin{important}[Parameter count]
		T0 theory works with the single parameter $\xi$ (plus the integer settings named in the corpus; SI/SM anchors such as $m_e$ serve for conversion):
		\begin{itemize}
			\item Standard Model of particle physics: 19+ parameters
			\item $\Lambda$CDM cosmology: 6 parameters
			\item T0 theory: 1 parameter ($\xi$)
		\end{itemize}
		This corresponds to a 96\% reduction in the number of parameters; the comparison with $\Lambda$CDM belongs to the cosmic sector and is speculative \textbf{[S]} (P39).
	\end{important}
	
	\section{Unit Analysis and Dimensional Consistency}
	
	\subsection{Verification of the Natural Units Framework}
	
	All T0 theory equations are dimensionally consistent in natural units:
	
	\begin{table}[h]
		\centering
		\adjustbox{max width=\textwidth}{%
\begin{tabular}{l l l l}
			\toprule
			Quantity & Natural units & Dimension & Verification \\
			\midrule
			$\xi$ & dimensionless & $[1]$ & satisfied \\
			$E_\xi$ & 7500 & $[E]$ & satisfied \\
			$L_\xi$ & $1.33 \times 10^{-4}$ & $[E^{-1}]$ & satisfied \\
			$T_\xi$ & 7500 & $[E]$ & satisfied \\
			$G_{\text{nat}}$ & $2.61 \times 10^{-70}$ & $[E^{-2}]$ & satisfied \\
			\bottomrule
		\end{tabular}%
}
		\caption{Dimensional consistency in natural units}
	\end{table}
	
	\subsection{Energy Scale Hierarchies}
	
	With the Planck energy as reference:
	
	\begin{align}
		E_{\text{Planck}} &= 1 \quad \text{(by definition in natural units)} \\
		E_\xi &= \frac{1}{\xi} = 7500
	\end{align}
	The simple powers $\xi^{1/2}\,E_{\text{Planck}} \approx 0.0115\,E_P = 1.4 \times 10^{17}$~GeV and $\xi^{1/3}\,E_{\text{Planck}} \approx 0.0511\,E_P = 6.2 \times 10^{17}$~GeV (with $E_{\text{Planck}} = 1.22 \times 10^{19}$~GeV), by contrast, do not give the weak and the QCD scale: the measured values are $v/E_P = 246/1.22 \times 10^{19} = 2.0 \times 10^{-17}$ (in the corpus: $v/E_P = \xi^4/(5\pi)$, Doc.~149) and $\Lambda_{\text{QCD}}/E_P \approx 0.2/1.22 \times 10^{19} = 1.6 \times 10^{-20}$; the powers $\xi^{1/2}$ and $\xi^{1/3}$ miss these scales by 15 and 18 orders of magnitude, respectively.
	
	\subsection{Additional Experimental Predictions}
	
	\textbf{Prediction 1: Electromagnetic resonance at characteristic $\xi$ frequency}
	\begin{itemize}
		\item Maximum $\xi$ field-photon coupling at $\nu = E_\xi = 7500$ (nat. units)
		\item Anomalies in electromagnetic propagation at this frequency
		\item Spectral features in the corresponding frequency range
	\end{itemize}
	
	\textbf{Prediction 2: Casimir force anomalies at characteristic $\xi$ length scale}
	\begin{itemize}
		\item Standard Casimir law: $F \propto d^{-4}$
		\item $\xi$ field modifications at $d \approx L_\xi = 10^{-4}$ m
		\item Measurable deviations through $\xi$ vacuum coupling
	\end{itemize}
	
	\textbf{Prediction 3: Modified vacuum fluctuations}
	\begin{itemize}
		\item Vacuum energy density variations at scale $L_\xi$
		\item Correlation between Casimir and CMB measurements
		\item Testable in precision laboratory experiments
	\end{itemize}
	
	\section{The Static Universe Paradigm}
	
	\subsection{Fundamental Properties of the T0 Universe}
	
	\begin{important}[Cosmological picture]
		\textbf{[S]} The T0 universe differs from expansion cosmology in the following points (speculative, P39):
		\begin{itemize}
			\item The universe does not expand
			\item The universe has existed eternally
			\item The universe has no beginning (no Big Bang)
			\item The universe is in thermodynamic equilibrium
			\item Cosmic phenomena are traced back to $\xi$ field dynamics
		\end{itemize}
	\end{important}
	
	\subsection{$r_0$ Definition from $\xi$}
	
	The fundamental length scale $r_0$ is defined by:
	\begin{align}
		r_0 &= \xi \cdot l_P = \frac{4}{3} \times 10^{-4} \times 1.616 \times 10^{-35}\,\text{m} \\
		&= 2.15 \times 10^{-39}\,\text{m}
	\end{align}
	
	In natural units with $l_P = 1$:
	\begin{equation}
		r_0 = \xi = \frac{4}{3} \times 10^{-4}
	\end{equation}
	
	\section{The Central Thesis: The Vacuum Is the $\xi$ Field}
	
	\begin{formula}
		The universal $\xi$ constant leads to the following connected relations:
		\begin{align}
			\xi &= \frac{4}{3} \times 10^{-4} \quad \text{(from geometry)} \\
			G &= \frac{\xi^2}{4m} \quad \text{(gravitation calculable)} \\
			T_{\text{CMB}} &= \frac{16}{9} \xi^2 \times E_\xi \quad \text{(unit-dependent match to about 1\,\%, see above)} \\
			\frac{|\rho_{\text{Casimir}}|}{\rho_{\text{CMB}}} &= \frac{\pi^2 \times 10^4}{960} \quad \text{(Casimir connection)}
		\end{align}
	\end{formula}
	
	\subsection{The Vacuum Is the $\xi$ Field}
	
	\begin{important}
		Central thesis of T0 theory:
		\begin{itemize}
			\item The vacuum is identical with the $\xi$ field
			\item The CMB is radiation of this vacuum at characteristic temperature \textbf{[S]}
			\item The Casimir force arises from geometric confinement of the same vacuum
			\item Gravitation follows from $\xi$ geometry
			\item The fundamental forces are traced back to $\xi$ field manifestations
		\end{itemize}
	\end{important}
	
	\subsection{Mathematical Structure}
	
	T0 theory formulates:
	\begin{enumerate}
		\item \textbf{Universal $\xi$ scaling}: The phenomena treated are derived from $\xi = \frac{4}{3} \times 10^{-4}$
		\item \textbf{Static paradigm} \textbf{[S]}: No Big Bang, no expansion, eternal existence
		\item \textbf{Time-energy consistency}: $T\cdot E = 1$ with lower time bound $T_0 = \xi\,t_P$
		\item \textbf{Dimensional consistency}: Completely formulated in natural units
		\item \textbf{Unit-independent physics}: Exact mathematical ratios
	\end{enumerate}
	
	\section{Conclusions}
	
	The T0 analysis of temperature units in natural units with CMB calculations yields:
	
	\begin{enumerate}
		\item \textbf{Universal $\xi$ scaling}: The temperature and energy scales treated follow from the geometric constant $\xi = \frac{4}{3} \times 10^{-4}$.
		
		\item \textbf{CMB without inflation} \textbf{[S]}: The theory interprets the CMB without inflation and without a decoupling at any cosmological $z$ (stationary $\xi$-field thermalization in an infinitely old universe), and derives primordial perturbations from T-field quantum fluctuations.
		
		\item \textbf{Static universe paradigm} \textbf{[S]}: In this picture, the universe is eternal and static and compatible with quantum mechanics without paradoxes.
		
		\item \textbf{Time-energy consistency} \textbf{[S]}: From $T\cdot E = 1$ with the lower time bound $T_0 = \xi\,t_P$ and the boundaryless torus it follows that there is no singular beginning (Doc.~306); the static universe does not require a Big Bang.
		
		\item \textbf{Dimensional consistency}: Consistent dimensions throughout in natural units, with $\xi$ as the only parameter (plus the geometric settings named).
		
		\item \textbf{Unit-independent physics}: All relations consist of exact mathematical ratios derived from fundamental geometry.
		
		\item \textbf{Testable predictions}: Specific, measurable deviations from $\Lambda$CDM that can be tested with next-generation experiments.
	\end{enumerate}
	
	\begin{important}[Conclusion]
		T0 theory proposes an alternative to expansion-based cosmology formulated in natural units \textbf{[S]} and describes temperature phenomena from particle physics to the cosmos with a single constant derived from geometry. The CMB calculations show that cosmological observations can be described within this framework; the cosmological part remains speculative (P39).
	\end{important}
	

	
	\begin{thebibliography}{20}
		\bibitem{T0Theory}
		Johann Pascher.
		\textit{The T0 Model (Planck-referenced): A Reformulation of Physics}.
		GitHub Repository, 2024.
		\url{https://jpascher.github.io/T0-Time-Mass-Duality/2/pdf}
		
		\bibitem{FineStructure}
		Johann Pascher.
		\textit{The Fine-Structure Constant: Different Representations and Relations}.
		Explains the critical distinction between $\alpha_{\text{EM}} = 1/137$ (SI) and $\alpha_{\text{EM}} = 1$ (natural units).
		2025.
		
		\bibitem{planck2020}
		Planck Collaboration (2020). 
		\textit{Planck 2018 Results. VI. Cosmological Parameters}. 
		Astronomy \& Astrophysics, 641, A6. 
		\url{https://doi.org/10.1051/0004-6361/201833910}
		
		\bibitem{codata2018}
		CODATA (2018). 
		\textit{The 2018 CODATA Recommended Values of the Fundamental Physical Constants}. 
		National Institute of Standards and Technology. 
		\url{https://physics.nist.gov/cuu/Constants/}
		
		\bibitem{casimir1948}
		Casimir, H. B. G. (1948). 
		\textit{On the Attraction Between Two Perfectly Conducting Plates}. 
		Proceedings of the Royal Netherlands Academy of Arts and Sciences, 51(7), 793--795.
		
		\bibitem{muon_g2_2021}
		Muon g-2 Collaboration (2021). 
		\textit{Measurement of the Positive Muon Anomalous Magnetic Moment to 0.46 ppm}. 
		Physical Review Letters, 126(14), 141801. 
		\url{https://doi.org/10.1103/PhysRevLett.126.141801}
		
		\bibitem{riess2022}
		Riess, A. G., et al. (2022). 
		\textit{A Comprehensive Measurement of the Local Value of the Hubble Constant with 1 km s$^{-1}$ Mpc$^{-1}$ Uncertainty from the Hubble Space Telescope and the SH0ES Team}. 
		The Astrophysical Journal Letters, 934(1), L7. 
		\url{https://doi.org/10.3847/2041-8213/ac5c5b}
		
		\bibitem{jwst_early}
		Naidu, R. P., et al. (2022). 
		\textit{Two Remarkably Luminous Galaxy Candidates at z $\approx$ 11--13 Revealed by JWST}. 
		The Astrophysical Journal Letters, 940(1), L14. 
		\url{https://doi.org/10.3847/2041-8213/ac9b22}
		
		\bibitem{cobe1992}
		COBE Collaboration (1992). 
		\textit{Structure in the COBE Differential Microwave Radiometer First-Year Maps}. 
		The Astrophysical Journal Letters, 396, L1--L5. 
		\url{https://doi.org/10.1086/186504}
	\end{thebibliography}
